\honest{One Life Against the World}{Eliezer Yudkowsky}{2007}

I open with the Talmud: ``Whoever saves a single life, it is as if he had saved the whole world.'' \nb{The line is from the Mishnah, Sanhedrin 4:5, where it is part of a court's warning to witnesses in capital cases. The printed Mishnah and the Babylonian Talmud read ``one soul from Israel''; the Jerusalem Talmud and the Kaufmann manuscript read ``one soul,'' as the post does.} It is a beautiful thought, and I ask you to feel its warm glow.

I can testify, I say, that helping one person feels just as good as helping the whole world. Once, when I was burned out and wasting time online, an anonymous blog comment of mine turned someone's whole life around. The high felt as good as a major scientific insight. \nb{The anecdote compares helping one person with an insight. The comparison with saving the world is the one the next sentences make with ``probably.''} Saving one life, I conclude, ``probably \emph{does} feel just as good as saving the entire world.''

But if you ever have to choose, save the world. Please. Beyond the warm glow is a gigantic difference.

Why might that not be obvious? I give three reasons: a qualitative duty to save what lives you can, which the person who saves one life and the person who saves the world both fulfill; the Greek ethics of personal virtue; and a value of one life so great that passing to the whole world changes little. \nb{The first view has a published defender. John Taurek argued in 1977 that no bystander is required to save five people rather than one, and proposed flipping a coin.}

I agree that one life is unimaginably valuable, and I hold that two lives are twice as unimaginably valuable. Whoever saves ten lives saves ten worlds, and whoever really saves the world saves an intergalactic civilization.

Then my argument. Two deaf children sleep on a railroad track. I drag one off, then stop and sip a Diet Pepsi while the train comes for the other, since I am already ``unimaginably'' far ahead on points. ``Doesn't sound right, does it?'' \nb{The second child can be saved at no cost to the first, so the first two views condemn this rescuer too. The case tells against the third view, read literally.}

Why should a philanthropist be any different, spending \$10 million on a cure for a disease that affects a hundred people when the same money could cure one that kills 10,000? ``I don't think it \emph{is} different.'' \nb{On the tracks both children could be saved; the philanthropist must choose. Views like Taurek's differ from the post only in such choices, and the post gives no further argument for this one.} When lives are at stake we have a duty to maximize, as strong as the duty to save lives. ``Whoever knowingly chooses to save one life, when they could have saved two \ldots\ they have damned themselves as thoroughly as any murderer.'' \nb{This equates letting die with killing. The Stanford Encyclopedia of Philosophy reports that consequentialists hold that view and anti-consequentialists ``almost universally'' reject it. The post gives no argument for it.}

In an addendum I say that obvious ways of spending money to save lives do not work or backfire, and I link a blog post and a magazine interview. \nb{The blog post is Robin Hanson's on the RAND Health Insurance Experiment, about free medical care. The interview is with James Shikwati, who argues that aid to Africa does more harm than good, a view others dispute.} Stuart Armstrong adds that if we disdain the philanthropist who spends inefficiently, we should disdain even more those who spend nothing.
